\documentclass[11pt]{article}
\usepackage[margin=1in]{geometry}
\usepackage{amsmath,amssymb}
\usepackage[hidelinks]{hyperref}

\newcommand{\Z}{\mathbb Z}
\newcommand{\R}{\mathbb R}

\title{A Later Small-Radii Literature Answer for the Discrete Euclidean Ball Maximal Problem}
\author{Run fa\_banach\_001}
\date{2026-07-19}

\begin{document}
\maketitle

\paragraph{Status.}
This is a lightweight literature-implied status note:
\[
\texttt{literature\_implied\_answer (partial small-radii subrange).}
\]
It records later progress on the discrete Euclidean ball maximal-function
problem highlighted in Bourgain--Mirek--Stein--Wrobel,
\emph{On the Hardy--Littlewood maximal functions in high dimensions:
Continuous and discrete perspective}, arXiv:1812.00153.

\paragraph{Source problem.}
In the introduction and in Section 3, ``Discrete perspective'', the source
studies whether the discrete Hardy--Littlewood maximal constant
\[
\mathcal C_p(d,G)
\]
for averages over \(G_t\cap\Z^d\) can be bounded independently of \(d\).
For \(G=B^2\), Theorem 1.2 of the source proves a dimension-free large-scale
estimate for
\[
\sup_{t\ge Cd}|\mathcal M_t^{B^2}f|.
\]
The source then says that the matter is reduced to understanding
\(\sup_{1\le N\le Cd}|\mathcal M_N^{B^2}f|\).  It records the dyadic ball
estimate for \(p\in[2,\infty]\), but states that the available multiplier
bounds do not give the full maximal function; in particular, obtaining the
third discrete multiplier estimate in equation (49) would be a route toward
dimension-free estimates for \(\mathcal C_p(d,B^2)\).

\paragraph{Supporting literature.}
Niksinski--Wrobel,
\emph{Dimension-free estimates for discrete maximal functions and lattice
points in high-dimensional spheres and balls with small radii},
arXiv:2503.16952, gives a later full-maximal small-radii result for the
Euclidean ball.  Its Theorem 1.1 states that for every \(\varepsilon>0\) there
is \(C_\varepsilon\) such that, uniformly in \(d\) and \(p\in[2,\infty]\),
\[
\left\|\sup_{0\le t\le d^{1/2-\varepsilon}}
|\mathcal M_t f|\right\|_{\ell^p(\Z^d)}
\le C_\varepsilon\|f\|_{\ell^p(\Z^d)}.
\]
This is a direct positive answer on a non-dyadic small-radii subrange of the
source problem.  Remark 1 of the supporting paper explicitly combines this
small-radii theorem with the large-scale theorem from the Bourgain--Mirek--
Stein--Wrobel survey to say that the remaining range relevant to Stein's
discrete question is
\[
t\in(d^{1/2-\varepsilon},Cd).
\]

\paragraph{Scope limitations.}
This packet is not a new proof and not a full settlement.  It records a later
partial literature answer.  The middle range
\[
d^{1/2-\varepsilon}<t<Cd
\]
remains outside this identification, as does removal of the \(\varepsilon\)
loss.  The supporting theorem is also in the range \(p\in[2,\infty]\), so it
does not settle the broader \(p\in(3/2,\infty]\) or \(p\in(1,\infty]\)
possibilities discussed in the source.

\paragraph{Search evidence.}
Cheap local indexes had no exact prior row for arXiv:1812.00153.  A web/arXiv
search for the exact paper and the discrete Euclidean-ball full maximal problem
identified arXiv:2503.16952 as the later small-radii result.  The parsed local
source and downloaded supporting PDF were checked for the source's Section 3
reduction and for Theorem 1.1 plus Remark 1 in the supporting paper.

\paragraph{References.}
\begin{itemize}
\item J. Bourgain, M. Mirek, E. M. Stein, and B. Wrobel, \emph{On the
Hardy--Littlewood maximal functions in high dimensions: Continuous and discrete
perspective}, arXiv:1812.00153, v2, 2019.
\item J. Niksinski and B. Wrobel, \emph{Dimension-free estimates for discrete
maximal functions and lattice points in high-dimensional spheres and balls with
small radii}, arXiv:2503.16952, v2, 2026.
\end{itemize}

\end{document}

